\section{模型检验与灵敏度分析}
\subsection{模型检验}
模型检验分别围绕时间对应、缺失预测与解释忠实性展开。

问题一中，100条样本均形成输出，1920个已生成词区间全部满足非负起点、正长度、音轨范围内终点与起点随词序单调四项约束。典型样本展示了词元、声学窗口与视频帧的对应关系，6条超长样本保留完整词级记录，定长接口截断不影响样本覆盖。

问题二中，验证集Macro-F1为0.6163、强度MAE为0.5818。273条分类错误中，213条涉及中性类，占78.0\%。测试集使用共用掩码，将七类输入处理与结构基线和两个公开基线与本文模型置于同一协议下比较。经来源视频分组配对bootstrap及Holm校正，18项局部比较中，相对模态丢弃的局部F1差异以及相对均值填充、三态嵌入、LMF、MulT的局部MAE差异达到5\%显著性水平。

问题三中，模态Shapley分解的效率残差不超过$8.9\times10^{-16}$，满足效率性公理。在96条验证子集上，高分片段删除相对随机删除的AOPC差为0.1068，95\%区间为[0.0887,0.1259]，说明局部响应排序与预测变化相符。该检验评价解释对当前模型的忠实性，主模态对基线与随机种子的敏感性另见表\ref{q3:tab:stability}。

\subsection{灵敏度分析}
时间对齐模型的灵敏度主要来自词边界与聚合参数。所选12条样本子集中，扰动次数由50增至200后，视觉方差的排序相关为0.995，锚点相对排序较为稳定。边界扰动尺度由40\,ms改为20或80\,ms时，音视频质量分数的排序相关均在0.90以上，说明该质量评价在所测范围内具有较好的排序一致性。

局部缺失预测的灵敏度从模态组合、连续区间比例、位置和长度四个方面检验。三模态居中缺口的目标长度由1增至16个锚点时，F1由0.6047变为0.4632，MAE由0.6496变为0.7981。逐样本对照显示，三模态共缺20\%（居中缺口）时60条样本由正确变为错误、48条由错误变为正确，16个锚点缺口下缺失时的分类错误中有50.1\%在完整输入下本已错误。短样本按$L_i-1$截断，并另报告各条件均满足局部删除与保留观测要求的共同子集，区分缺失程度变化与样本边界效应。

解释模型的灵敏度表现为少数样本的主模态变化。B1与B2在验证子集上的$\kappa$为0.5686，B1与B3为$-0.0231$；B3下原有10条非文本主导样本均转为文本。文本在总体上占主导，而非文本贡献需要结合解释基线考察。

\section{模型评价与推广}
\subsection{模型特点}
\textbf{时间组织明确。} 以原词区间连接子词表示与音视频帧，用交叠权重统一不同采样粒度，并以观测掩码保留缺失状态，形成从定长特征到原始素材的对应关系。

\textbf{缺失处理与任务学习相结合。} 三态输入区分结构填充、有效观测和局部缺失，补全表示通过门控参与融合，EMT-DLFR教师提供分类与强度软监督。局部输入模态消融及连续缺口实验分别刻画模态敏感性与缺失程度的影响。

\textbf{解释对应同一预测器。} 模态联盟分解与局部删除共享预测模型，分别给出模态贡献和片段响应，再通过词级时间映射定位原始证据，使预测数值与解释结果相互对应。

\subsection{模型局限与改进}
时间对齐采用自动估计的词区间，当前验证侧重覆盖、结构与相对时序一致性。可进一步加入人工词边界标注，评价起止点误差，并改善数字读法与长序列处理。

缺失预测的分类误差集中于中性与非中性类别之间，逆频率加权均衡了训练目标中的类别贡献，但类别识别仍取决于样本的有效证据。后续可加强缺口上下文与弱情感样本的区分，通过组件消融和独立数据进一步检验补全、蒸馏的作用。

解释结果随基线和训练种子变化，非文本主导样本尤为明显。新增冻结测试子集分别确认了分类与强度片段的删除响应，双任务证据均已回映至原素材；进一步改进可结合独立数据与人工时间标注，检验解释迁移性及定位误差。

在应用推广中，词级锚点与掩码表示可用于其他具有转写文本的音视频分析任务；模型结构和参数需结合新任务的特征维度、缺失模式与标签分布重新确定。
